% ============================================================
% f2014_tianjin_science_q17_pyramid.tex
% 2014 天津卷（理科）第 17 题：四棱锥 P-ABCD
% 来源 PDF：普通高考/2014/2014天津理.pdf
% 原页：第 2 页（A3），印面 924，第 17 题题干下方
% PA⊥底面 ABCD，AD⊥AB，AB∥DC，AD=DC=AP=2，AB=1，E 为 PC 中点。
% 由三视图式平行投影绘制，不按严格比例；PB、BE 为可见棱，
% PD、DC、AD、AC、BD 为不可见线（虚线）。
% ============================================================
\begin{tikzpicture}[>=Stealth, line width=0.65pt, font=\small, scale=1.0,
  line cap=round, line join=round]
  \coordinate (A) at (0,0);
  \coordinate (B) at (2.00,0);
  \coordinate (D) at (1.76,1.47);
  \coordinate (C) at (5.76,1.47);
  \coordinate (P) at (0,3.94);
  \coordinate (E) at (2.88,2.705);

  % 可见棱
  \draw (P) -- (A);
  \draw (P) -- (B);
  \draw (P) -- (C);
  \draw (A) -- (B);
  \draw (B) -- (C);
  \draw (B) -- (E);

  % 不可见线
  \draw[dashed] (P) -- (D);
  \draw[dashed] (A) -- (D);
  \draw[dashed] (D) -- (C);
  \draw[dashed] (A) -- (C);
  \draw[dashed] (B) -- (D);

  \fill (A) circle (0.85pt) node[below left=-1pt] {$A$};
  \fill (B) circle (0.85pt) node[below=-1pt] {$B$};
  \fill (C) circle (0.85pt) node[right=-1pt] {$C$};
  \fill (D) circle (0.85pt) node[above left=-1pt] {$D$};
  \fill (P) circle (0.85pt) node[above=-1pt] {$P$};
  \fill (E) circle (0.85pt) node[above right=-2pt] {$E$};
\end{tikzpicture}
